\documentclass[troisieme]{fiche}
\metadonnees{17}{Il n'y avait pas de garçon}{Bloc D — Rapporter et supposer}

\begin{document}
\entete

\objectifs{%
\textbf{Langue} : le discours indirect ; \emph{say}, \emph{tell}, \emph{ask}.\par
\textbf{Texte} : L.\,M. Montgomery, \emph{Anne of Green Gables} (1908), ch. \textsc{iii}.\par
\textbf{Durée} : 1 h — exercices 20 min, texte et questions 25 min, version 15 min.}

%% ===================================================================
\exercice[12 min]{Rapporter des paroles}

\rappel{\textbf{Rappel.} Quand le verbe introducteur est au passé, le temps rapporté recule
d'un cran : présent $\rightarrow$ prétérit, \emph{will} $\rightarrow$ \emph{would}.}

\consigne{Rapportez les paroles suivantes.}

\begin{anglais}
\begin{enumerate}[label=\arabic*.,itemsep=3mm,leftmargin=*]
  \item ``I am an orphan,'' she said.
    \lignes{1}
    \corr{She said (that) she was an orphan.}
  \item ``We want a boy,'' said Marilla.
    \lignes{1}
    \corr{Marilla said (that) they wanted a boy.}
  \item ``Where is the boy?'' she asked.
    \lignes{1}
    \corr{She asked where the boy was.}
  \item ``Call me Cordelia,'' said the child.
    \lignes{1}
    \corr{The child asked Marilla to call her Cordelia.}
  \item ``I can't sleep,'' said Anne.
    \lignes{1}
    \corr{Anne said (that) she couldn't sleep.}
  \item ``Do you like the name?'' he asked me.
    \lignes{1}
    \corr{He asked me if I liked the name.}
  \item ``Don't cry,'' Marilla told her.
    \lignes{1}
    \corr{Marilla told her not to cry.}
  \item ``I will stay here,'' she said.
    \lignes{1}
    \corr{She said (that) she would stay there.}
  \item ``Mrs Spencer is making a mistake,'' said Matthew.
    \lignes{1}
    \corr{Matthew said (that) Mrs Spencer was making a mistake.}
  \item ``Are you ashamed of your name?'' Marilla asked the child.
    \lignes{1}
    \corr{Marilla asked the child if she was ashamed of her name.}
\end{enumerate}
\end{anglais}

\note[Trois choses changent]{Le temps recule (1, 2, 5, 8, 9). Les pronoms suivent celui qui
rapporte : \emph{me} devient \emph{her} (4). Les mots qui situent changent aussi :
\emph{here} devient \emph{there} (8), \emph{tomorrow} devient \emph{the next day}.}

\note[Questions et ordres]{Une question rapportée redevient une phrase ordinaire : plus
d'inversion, plus de \emph{do} --- \emph{she asked where the boy was} (3, 6, 10). Une question
sans mot interrogatif prend \emph{if} ou \emph{whether}. Un ordre devient \emph{tell} ou
\emph{ask} + \emph{to} + base verbale (4, 7).}

%% ===================================================================
\exercice[8 min]{\emph{Say}, \emph{tell}, \emph{ask}}

\rappel{\textbf{Rappel.} \emph{Tell} se construit avec la personne : \emph{tell me}.
\emph{Say} non : \emph{say to me} ou, mieux, \emph{say} tout court.}

\consigne{Traduisez en anglais.}

\begin{anglais}
\begin{enumerate}[label=\arabic*.,itemsep=3mm,leftmargin=*]
  \item Il m'a dit qu'il était fatigué.
    \lignes{1}
    \corr{He told me (that) he was tired.}
  \item Elle a dit qu'elle ne viendrait pas.
    \lignes{1}
    \corr{She said (that) she would not come.}
  \item Je lui ai demandé où il habitait.
    \lignes{1}
    \corr{I asked him where he lived.}
  \item Il m'a demandé de partir.
    \lignes{1}
    \corr{He asked me to leave.}
  \item Elle nous a dit de ne pas faire de bruit.
    \lignes{1}
    \corr{She told us not to make any noise.}
  \item Demande-lui s'il a compris.
    \lignes{1}
    \corr{Ask him if he has understood.}
\end{enumerate}
\end{anglais}

\note[Le partage]{\emph{Tell} + personne + ce qu'on dit (1, 5). \emph{Say} + ce qu'on dit,
sans personne (2). \emph{Ask} pour une question (3, 6) comme pour une demande (4). \emph{He
said me} est une faute ; \emph{he told that} aussi.}

%% ===================================================================
\newpage
\section{Texte}

\noindent\small\textit{Matthew et Marilla Cuthbert, frère et sœur, ont demandé un garçon à
l'orphelinat pour les aider à la ferme. Matthew revient de la gare avec une fillette de onze
ans.}\normalsize

\begin{texte}
Marilla came \gl[briskly]{briskly}{d'un pas vif} forward as Matthew opened the door. But when
her eyes fell on the odd little figure in the stiff, ugly dress, with the long
\gl[a braid]{braids}{une natte} of red hair and the eager, \gl[luminous]{luminous}{lumineux}
eyes, she stopped short in \gl[amazement]{amazement}{stupéfaction}.

``Matthew Cuthbert, who's that?'' she \gl[to ejaculate]{ejaculated}{lâcher, s'exclamer}.
``Where is the boy?''

``There wasn't any boy,'' said Matthew \gl[wretchedly]{wretchedly}{piteusement}. ``There was
only \emph{her}.'' He nodded at the child, remembering that he had never even asked her name.

``No boy! But there \emph{must} have been a boy,'' insisted Marilla. ``We sent word to Mrs.
Spencer to bring a boy.''

``Well, she didn't. She brought \emph{her}. I asked the stationmaster. And I had to bring her
home. She couldn't be left there, no matter where the mistake had come in.''

``Well, this is a pretty piece of business!'' ejaculated Marilla.

During this dialogue the child had remained silent, her eyes \gl[to rove]{roving}{errer} from
one to the other, all the animation fading out of her face. Suddenly she seemed to
\gl[to grasp]{grasp}{saisir} the full meaning of what had been said. Dropping her precious
carpet-bag she sprang forward a step and \gl[to clasp]{clasped}{joindre, serrer} her hands.

``You don't want me!'' she cried. ``You don't want me because I'm not a boy! I might have
expected it. Nobody ever did want me. I might have known it was all too beautiful to last. Oh,
what shall I do? I'm going to \gl[to burst]{burst}{éclater} into tears!''

Burst into tears she did. Sitting down on a chair by the table, flinging her arms out upon it,
and burying her face in them, she proceeded to cry \gl[stormily]{stormily}{à gros sanglots}.
Marilla and Matthew looked at each other
\gl[deprecatingly]{deprecatingly}{d'un air désolé} across the stove. Neither of them knew what
to say or do. Finally Marilla stepped \gl[lamely]{lamely}{maladroitement} into the
\gl[a breach]{breach}{une brèche}.

``Well, well, there's no need to cry so about it.''

``Yes, there \emph{is} need!'' The child raised her head quickly, revealing a
\gl[tear-stained]{tear-stained}{barbouillé de larmes} face and trembling lips. ``\emph{You}
would cry, too, if you were an orphan and had come to a place you thought was going to be
home and found that they didn't want you because you weren't a boy. Oh, this is the most
\emph{tragical} thing that ever happened to me!''

Something like a \gl[reluctant]{reluctant}{réticent} smile, rather \gl[rusty]{rusty}{rouillé}
from long \gl[disuse]{disuse}{manque d'usage}, \gl[to mellow]{mellowed}{adoucir} Marilla's
\gl[grim]{grim}{sévère} expression.

``Well, don't cry any more. We're not going to turn you out-of-doors tonight. You'll have to
stay here until we investigate this affair. What's your name?''

The child hesitated for a moment.

``Will you please call me Cordelia?'' she said eagerly.
\end{texte}
\source{L.\,M. Montgomery, \emph{Anne of Green Gables}, Boston, Page, 1908,
ch.~\textsc{iii}.}

\partiel[15 min]{Questions}
\consigne{Répondez aux deux premières questions en français, aux deux dernières en anglais.}

\begin{enumerate}[label=\arabic*.,itemsep=2.5mm,leftmargin=*]
  \item Pourquoi Marilla s'arrête-t-elle net en voyant l'enfant ?
    \lignes{2}
    \corr{Parce qu'elle attendait un garçon et découvre une petite fille : une silhouette
    étrange, une robe raide et laide, de longues nattes rousses et de grands yeux brillants.}
  \item Comment Matthew explique-t-il ce qui s'est passé ?
    \lignes{3}
    \corr{Il n'y avait pas de garçon : Mrs Spencer a amené cette enfant. Il s'est renseigné
    auprès du chef de gare et a dû la ramener, car on ne pouvait pas l'abandonner là, d'où que
    vienne l'erreur.}
  \item \en{``I might have expected it. Nobody ever did want me.'' What do these two sentences
    tell us about the child's life so far?}
    \lignes{3}
    \corr{That being rejected is what she is used to: she blames herself for having hoped, not
    the adults for sending her back. \emph{Nobody ever did want me} — the \emph{did} makes the
    sentence heavier, as if she had been told so many times.}
  \item \en{Find in the text one sentence of reported speech and turn it back into direct
    speech.}
    \lignes{3}
    \corr{\emph{We sent word to Mrs. Spencer to bring a boy} $\rightarrow$ \og Bring us a
    boy \fg. \emph{He had never even asked her name} $\rightarrow$ \og What is your name?
    \fg{} qu'il n'a pas posée. \emph{What had been said} reprend tout le dialogue sans le
    répéter : c'est le résumé le plus court possible, et c'est lui qui frappe l'enfant.}
\end{enumerate}

\note[Ce que le texte ne dit pas]{Personne ne s'adresse à l'enfant pendant toute la première
moitié de la scène. On parle d'elle à la troisième personne devant elle — \emph{there was only
her} —, et c'est ce silence qu'elle finit par comprendre.}

%% ===================================================================
\section{Version}
\consigne{Traduisez le passage en français. Il suit immédiatement le texte.}

\begin{version}
``\emph{Call} you Cordelia? Is that your name?''

``No-o-o, it's not exactly my name, but I would love to be called Cordelia. It's such a
perfectly elegant name.''

``I don't know what on earth you mean. If Cordelia isn't your name, what is?''

``Anne Shirley,'' reluctantly \gl[to falter]{faltered}{balbutier} forth the owner of that
name, ``but, oh, please do call me Cordelia. It can't matter much to you what you call me if
I'm only going to be here a little while, can it? And Anne is such an unromantic name.''

``Unromantic \gl[fiddlesticks]{fiddlesticks}{balivernes}!'' said the unsympathetic Marilla.
``Anne is a real good \gl[plain]{plain}{simple, ordinaire} sensible name. You've no need to be
ashamed of it.''
\end{version}
\source{\emph{Ibid.}}

\lignes{16}

\corr{\textbf{Proposition de traduction.} \og Vous \emph{appeler} Cordelia ? Est-ce là votre
nom ? \fg{} \og No-o-on, ce n'est pas exactement mon nom, mais j'adorerais qu'on m'appelle
Cordelia. C'est un nom d'une élégance parfaite. \fg{} \og Je ne comprends pas un mot de ce que
vous racontez. Si Cordelia n'est pas votre nom, quel est-il ? \fg{} \og Anne Shirley \fg,
balbutia à contrecœur celle qui le portait, \og mais, oh, appelez-moi Cordelia, je vous en
prie. Cela ne peut pas faire grande différence pour vous, si je ne dois rester ici qu'un petit
moment, n'est-ce pas ? Et Anne est un nom si peu romanesque. \fg{} \og Peu romanesque,
balivernes ! dit Marilla sans s'attendrir. Anne est un nom simple, solide et raisonnable. Vous
n'avez pas à en avoir honte. \fg}

\note[Choix de traduction 1]{\emph{I would love to be called Cordelia} : passif à l'infinitif.
Le français passe par \og que l'on m'appelle \fg{} — il n'a pas d'infinitif passif léger.}

\note[Choix de traduction 2]{\emph{sensible} est un faux ami : \og raisonnable, sensé \fg.
\og Sensible \fg{} au sens français se dit \emph{sensitive}. La phrase de Marilla oppose le
solide au romanesque : tout le contresens tiendrait à ce mot.}

\note[Choix de traduction 3]{\emph{faltered forth the owner of that name} : la narratrice
évite de la nommer au moment même où elle se nomme. Garder le détour --- \og celle qui le
portait \fg{} --- et non \og Anne \fg.}

%% ===================================================================
\partiel{Lexique à mémoriser}
\begin{multicols}{2}\small
\begin{itemize}[label=--,itemsep=0.8mm,leftmargin=*]
  \item \textbf{an orphan} — un orphelin
  \item \textbf{a tear} — une larme
  \item \textbf{to cry} — pleurer ; \textit{et} crier
  \item \textbf{ashamed} — honteux ; \textit{to be ashamed of} avoir honte de
  \item \textbf{to spell} — épeler
  \item \textbf{plain} — simple, ordinaire
  \item \textbf{sensible} — raisonnable \textit{(faux ami)}
  \item \textbf{to hesitate} — hésiter
  \item \textbf{to nod} — hocher la tête
  \item \textbf{a mistake} — une erreur \textit{(fiche 16)}
  \item \textbf{eagerly} — avec empressement
  \item \textbf{to insist} — insister
\end{itemize}
\end{multicols}

\end{document}
